\section{Discussão}
\label{sec:discussao}

\subsection{Lacunas do estado da arte}

A comparação dos \nIncluidos{} estudos analisados evidencia cinco lacunas. São
lacunas \emph{nos estudos analisados}: os estudos analisados vêm de uma fase
exploratória e não de um critério de seleção definido \emph{a priori}
(Secção~\ref{sec:selecao}), pelo que não se afirma que nenhum dos
restantes estudos elegíveis as resolva.

\begin{enumerate}[label=\textbf{L\arabic*.}, leftmargin=*]
  \item \textbf{A frescura e a deteção de duplicados não aparecem combinadas.}
    Malina et al.\ e Van Glabbeek et al.\ assinam um carimbo temporal e
    apresentam-no como proteção contra repetição, mas não definem a janela de
    aceitação nem eliminam duplicados dentro dela~\cite[§5.2]{malina2019secure},
    \cite[§3.4]{vanglabbeek2022flexible}. O ENTRUST deteta duplicados com um
    contador por publicador, mas o seu campo temporal indica quando chega a
    próxima mensagem, e não a idade máxima de cada
    mensagem~\cite[§5.3, §6]{dahlmanns2021entrust}.
  \item \textbf{A ligação da mensagem ao contexto é incompleta.} Malina et al.\
    assinam o tópico e o nível de segurança~\cite[§4.5]{malina2019secure}, e o
    ENTRUST usa uma chave por tópico~\cite[§5]{dahlmanns2021entrust}. Nenhum dos
    estudos analisados aplica, no recetor, uma restrição entre a credencial
    autenticada e os tópicos em que pode publicar, nem uma política local que
    recuse níveis de proteção mais fracos.
  \item \textbf{O \emph{broker} continua, em regra, a ser de confiança.} Só quatro
    dos 16 mecanismos dispensam, total ou parcialmente, essa confiança, e vários
    assumem-na explicitamente~\cite[§III-B]{singh2015secure},
    \cite[§IV-A]{calabretta2018mqttauth}, \cite[§4.6]{malina2019secure}.
  \item \textbf{O não repúdio é raro e depende de terceiros confiáveis.} Só Malina
    et al.\ e Van Glabbeek et al.\ o oferecem, com um \emph{broker} que decifra ou
    é honesto mas curioso e, no segundo caso, uma CA totalmente
    confiável~\cite[§3.1]{vanglabbeek2022flexible}. As restantes soluções usam
    segredos partilhados, pelo que quem os detém, incluindo o recetor, pode forjar
    mensagens em nome do remetente.
  \item \textbf{Custo e benefício não são medidos em conjunto.} As avaliações
    medem o custo, mas não o relacionam, no mesmo ensaio, com os ataques impedidos,
    e usam plataformas e métricas que não são comparáveis entre si.
\end{enumerate}

Nenhuma destas lacunas pede um mecanismo criptográfico novo: os elementos
existem, isolados, na literatura. A arquitetura proposta responde a L1, L2 e L4
combinando-os no mesmo recetor e sem confiar no \emph{broker} (L3), e o desenho em
camadas comutáveis foi pensado para que a segunda parte do projeto responda a L5
(Secção~\ref{sec:plano}).

\subsection{Ameaças à validade}

A Tabela~\ref{tab:validade} resume as ameaças à validade da revisão, seguindo
Kitchenham e Charters~\cite{kitchenham2007}.

\begin{table}[ht]
  \centering
  \caption{Ameaças à validade da revisão e mitigação.}
  \label{tab:validade}
  \small
  \begin{tabularx}{\textwidth}{@{}L{2.6cm}YY@{}}
    \toprule
    \textbf{Tipo} & \textbf{Ameaça} & \textbf{Mitigação} \\
    \midrule
    Construção & A \emph{string} pode não captar estudos que usem outra
      terminologia & Grupos de sinónimos; refinamento após a leitura das revisões;
      inclusão de MQTT-SN e de termos de deteção; os estudos já conhecidos foram
      todos devolvidos pelas pesquisas \\
    Seleção & Enviesamento na aplicação dos critérios na triagem, feita por um só
      revisor, sem medida de concordância & Critérios operacionalizados com exemplos;
      regra conservadora (na dúvida, o registo segue); decisão e critério de cada
      registo publicados no repositório, para verificação independente \\
    Síntese & Só \nIncluidos{} dos \nElegiveis{} estudos elegíveis foram analisados em
      texto integral, identificados numa fase exploratória e não escolhidos por um
      critério \emph{a priori}; as lacunas podem estar tratadas em estudos não
      analisados & Desvio declarado; todos confirmados como elegíveis; lacunas
      formuladas \enquote{nos estudos analisados}; lista completa da triagem no
      repositório~\cite{mqttsecure2026repo} \\
    Extração & Subjetividade na classificação das propriedades; um estudo pode
      garantir uma propriedade sem a descrever & Leitura do texto integral;
      \enquote{não reportado} distinto de \enquote{ausente por desenho}; secção e
      página de cada célula registadas no repositório \\
    Externa & Só três bases de dados; ACM DL e Web of Science não pesquisadas;
      enviesamento de publicação & O Scopus indexa as atas ACM relevantes (p.\ ex.\
      ARES, SaT-CPS, ICPE); inglês e português; revista, conferência e
      \emph{workshop} \\
    Conclusão & Os números por propriedade dependem da granularidade das
      categorias & Taxonomia e critérios publicados; dados de extração no
      Anexo~\ref{anx:extracao} \\
    \bottomrule
  \end{tabularx}
\end{table}

A ameaça mais relevante é a da síntese: com \nElegiveis{} estudos elegíveis e
\nIncluidos{} analisados, as percentagens da Secção~\ref{sec:resultados} descrevem o
conjunto analisado, e não todo o campo. Na extração, a leitura integral alterou
várias células face à leitura dos resumos (por exemplo, o contador anti-repetição
do ENTRUST e o não repúdio de Malina et al.), o que confirma que \enquote{não
reportado no resumo} não deve ser lido como \enquote{ausente}.
